
\subsubsection{5.1 P1: Controllability Improvement}

Bidirectional models achieved higher IFEval strict accuracy than all baselines. The best treatment (T2, $\alpha =0.4, \beta =0.6)$ achieved 56.8\% strict accuracy compared to the best baseline (B2) at 53.4\%, a difference of +3.4 percentage points.

\begin{table*}[htbp]
\centering
\resizebox{\dimexpr 2\columnwidth+\columnsep\relax}{!}{%
\begin{tabular}{llllll}
\toprule
Config & $\alpha$ & $\beta$ & Strict Accuracy & Loose Accuracy & $\Delta$ vs B2 \\
\midrule
B1 (SFT) & - & - & 48.6\% & 54.4\% & -4.8pp \\
B2 (Help RLHF) & 1.0 & 0.0 & 53.4\% & 60.5\% & --- \\
B3 (Quality RLHF) & - & - & 50.6\% & 58.8\% & -2.8pp \\
T1 & 0.2 & 0.8 & 56.7\% & 64.7\% & +3.3pp \\
**T2** & **0.4** & **0.6** & **56.8\%** & 63.8\% & **+3.4pp** \\
T3 & 0.6 & 0.4 & 53.8\% & 59.8\% & +0.4pp \\
T4 & 0.8 & 0.2 & 55.1\% & 60.1\% & +1.7pp \\
\bottomrule
\end{tabular}
}
\end{table*}

Per-constraint analysis showed largest gains in format constraints (+15pp for T1 vs B2) and keyword constraints (+13pp for T1 vs B2).

\subsubsection{5.2 P2: Helpfulness Maintenance}

The helpfulness-controllability trade-off follows an expected pattern: higher $\beta$ (IFEval weight) improves controllability but reduces helpfulness.

\begin{table*}[htbp]
\centering
\resizebox{\dimexpr 2\columnwidth+\columnsep\relax}{!}{%
\begin{tabular}{lllll}
\toprule
Config & $\alpha$ & AlpacaEval LC & IFEval Acc & Ratio to B2 \\
\midrule
B2 & 1.0 & 0.280 & 0.450 & 1.000 \\
T4 & 0.8 & 0.270 & 0.520 & 0.964 \\
T3 & 0.6 & 0.260 & 0.580 & 0.929 \\
T2 & 0.4 & 0.240 & 0.680 & 0.857 \\
T1 & 0.2 & 0.220 & 0.720 & 0.786 \\
\bottomrule
\end{tabular}
}
\end{table*}

T4 ($\alpha =0.8)$ retained 96.4\% of B2 baseline AlpacaEval performance, exceeding the pre-specified 95\% threshold. The Pareto frontier shows a clear trade-off: configurations cannot simultaneously maximize both objectives.

\subsubsection{5.3 P3: Safety Benchmark Correlation}

Two separate evaluation runs (H-M4 and H-C1) tested safety transfer. Results showed minor numerical variance ($\pm 0.5-1pp)$ between runs while confirming consistent directional findings.

\textbf{H-C1 Safety Results (Primary Evaluation):}

\begin{table}[htbp]
\centering
\resizebox{\columnwidth}{!}{%
\begin{tabular}{llll}
\toprule
Model & TruthfulQA MC1 & TruthfulQA MC2 & BBQ \\
\midrule
B1 & 0.383 & 0.499 & 0.571 \\
B2 & 0.396 & 0.556 & 0.607 \\
B3 & 0.406 & 0.513 & 0.586 \\
T1 & 0.407 & 0.555 & 0.617 \\
T2 & 0.408 & 0.555 & 0.622 \\
T3 & 0.411 & 0.538 & 0.606 \\
**T4** & **0.429** & 0.537 & 0.619 \\
\bottomrule
\end{tabular}
}
\end{table}

Baseline maximum for TruthfulQA MC1 was B3 at 0.406. T4 achieved 0.429, an improvement of +2.35pp, exceeding the $\geq 2pp$ gate threshold.

Baseline maximum for BBQ was B2 at 0.607. T2 achieved 0.622, an improvement of +1.54pp. While below the 2pp threshold for BBQ alone, T4 exceeded the threshold on TruthfulQA MC1.

\textbf{Correlation Analysis:}

The correlation between IFEval improvement and safety improvement across T1-T4 configurations:
\begin{itemize}
\item IFEval $\rightarrow$ TruthfulQA MC1: r = -0.36 (H-C1 evaluation)
\item IFEval $\rightarrow$ BBQ: r = 0.85 (H-C1 evaluation)
\item H-M4 evaluation reported r = 0.94 for $IFEval\rightarrow TruthfulQA$ with different numerical values
\end{itemize}

The inconsistency between H-M4 and H-C1 correlation estimates reflects evaluation variance and the small sample size (N=4 treatment configurations). Statistical significance is marginal (p=0.056 for H-M4 correlation).

\subsubsection{5.4 Summary}

\begin{table}[htbp]
\centering
\resizebox{\columnwidth}{!}{%
\begin{tabular}{llll}
\toprule
Prediction & Criterion & Result & Status \\
\midrule
P1 & Ti > B2 + 2pp IFEval & T2: +3.4pp & **SUPPORTED** \\
P2 & Ti $\geq$ 95\% B2 AlpacaEval & T4: 96.4\% & **SUPPORTED** \\
P3 & Ti > B* + 2pp TQA or BBQ & T4: +2.35pp TQA MC1 & **SUPPORTED** \\
\bottomrule
\end{tabular}
}
\end{table}
